\begin{dang}{Nhiệt lượng cần cung cấp}
	\Anlythuyet{\setcounter{paragraph}{0}
		\paragraph{Nhiệt lượng cần cung cấp để làm nóng chảy hoàn toàn vật có khối lượng m}  
		\begin{align*}
			\boder{Q = Q_1+Q_2}\quad
			\begin{cases}
				Q_1\ \ct{là nhiệt lượng cần cung cấp để vật rắn tăng nhiệt độ lên tới}\\ \quad \ \   \ct{nhiệt độ nóng chảy}\\
				Q_2\ \ct{là nhiệt nóng chảy của vật}
			\end{cases}
		\end{align*}
		\paragraph{Nhiệt lượng cần cung cấp để hoá hơi hoàn toàn chất lỏng có khối lượng m} 
		\begin{align*}
			\boder{Q = Q_1+Q_2}\quad 
			\begin{cases}
				Q_1\ \ct{là nhiệt lượng cần cung cấp để chất lỏng tăng nhiệt độ lên tới}\\ 
				\quad \ \  \ct{nhiệt độ sôi}\\
				Q_2\ \ct{là nhiệt hoá hơi của chất lỏng}
			\end{cases}
	\end{align*}}
\end{dang}
